\begin{abstract}
% 基于可验证奖励的强化学习（RLVR）通过奖励最终答案的正确性提升了大语言模型的数学推理能力，但并未对中间推理过程进行监督。因此，模型可能通过缺乏依据的推导、循环依赖或不必要地冗长的推理轨迹得到正确答案。为此，我们提出 \textbf{TopoPRM}，一种可验证的过程奖励框架，在标准结果奖励之外引入拓扑感知监督。对于每条推理轨迹，我们使用基于规则的解析器构建步骤间依赖关系的有向无环图（DAG）。在此基础上，我们定义了两类确定性过程奖励：\emph{拓扑结构奖励}用于约束依赖图的无环性、连通性和方向一致性，\emph{连续性奖励}用于检查每一步是否建立在先前已确立的表达式或结论之上。这两类奖励均由确定性程序计算，无需训练额外的奖励模型，也不依赖人工标注。为防止过程奖励偏向那些结构上有效但结论错误的推理轨迹，我们进一步提出 \emph{Stratified Clipping Advantage Estimation}（SCAE），即在正确性分层内计算优势。基于一个在人工标注数学 critique 样本上训练得到的监督式冷启动模型，我们将基于 LoRA 的 GRPO 应用于数学推理任务，并在涵盖不同难度水平的公开基准上进行评测。相较于仅使用结果奖励的 GRPO，TopoPRM 在提升答案准确率的同时，还能生成更短且更连贯的推理轨迹。进一步地，reverse-KL 蒸馏阶段可将这些收益迁移到更小的学生模型上，同时降低推理成本。
Reinforcement learning with verifiable rewards (RLVR) is an effective
paradigm for improving mathematical reasoning in LLMs. Yet the optimization signal is typically dominated by final-answer correctness in the practical training pipeline, leaving the quality of intermediate reasoning steps only weakly constrained. 
Consequently, models may reach correct answers through unsupported deductions, circular dependencies, or unnecessarily long reasoning traces. 
We propose \textbf{TopoPRM}, a verifiable process-reward framework that introduces structural supervision without relying on learned reward models. 
A deterministic parser extracts an implicit directed acyclic graph (DAG) from each reasoning trace.
For each reasoning trace, a rule-based parser constructs a directed acyclic graph (DAG) over inter-step dependencies. 
We then compute two process signals: a\emph{topological structure reward} for acyclicity, conclusion support, directional consistency, and optional reference-edge coverage; and a \emph{continuity reward} that verifies step traceability to prior claims or explicit givens.
%Both rewards are computed by deterministic programs and require neither learned reward models nor human annotations. 
To combine these signals with outcome, format, and length rewards while preserving correctness primacy, we adopt \emph{Stratified Clipping Advantage Estimation} (SCAE) as an accuracy-first stratified shaping strategy in GRPO training.
Built on a supervised cold-start model trained on human-annotated mathematical critique samples, we apply LoRA-based GRPO to mathematical reasoning and evaluate on public benchmarks of varying difficulty. 
TopoPRM improves answer accuracy over outcome-only GRPO while producing shorter and more coherent reasoning traces. A reverse-KL distillation stage further transfers these gains to smaller student models at lower inference cost.
\end{abstract}

% Reinforcement learning with verifiable rewards (RLVR) improves final-answer accuracy in mathematical reasoning, yet the optimisation signal is typically dominated by outcome correctness, leaving intermediate reasoning quality only weakly constrained.
% We propose \textbf{TopoPRM}, a verifiable process-reward framework that introduces structural supervision without learned reward models.
% For each reasoning trace, a rule-based parser constructs a directed acyclic graph (DAG) over inter-step dependencies.
% We then compute two deterministic process signals: a \emph{topological structure reward} for acyclicity, conclusion support, and directional consistency; and a \emph{continuity reward} that verifies step traceability to prior claims or problem givens.
% To combine these signals with outcome, format, and length rewards while preserving correctness primacy, we adopt hierarchical reward aggregation with stratified clipping in GRPO training.
% A reverse-KL distillation stage further transfers these gains to smaller student models at lower inference cost.
% On private mathematical critique benchmarks, TopoPRM improves critique accuracy by $+13.4$ points over outcome-only GRPO while reducing response length by 11\%.
% The distilled Qwen3-8B student achieves 82.5\% on GSM8K and 59.8\% on MATH-500, competitive with models of comparable scale.